\section{Optimalizační modely pro podporu ekonomického rozhodování. Příklady typických úloh LP a MILP, heuristické optimalizační algoritmy}

\subsection{Základní myšlenka optimalizace}

Optimalizace slouží k tomu, abychom z množiny přípustných rozhodnutí vybrali takové, které je
nejlepší vzhledem ke zvolenému kritériu. V ekonomii a managementu jde typicky o maximalizaci
zisku, výnosu nebo užitku, případně minimalizaci nákladů, času, rizika, vzdálenosti či ztráty.

Obecný optimalizační problém lze zapsat jako
\[
\begin{aligned}
\min/\max \quad & f(x)\\
\text{s.t.}\quad & g_i(x)\le 0,\quad i=1,\dots,m,\\
& h_j(x)=0,\quad j=1,\dots,p,\\
& x\in X.
\end{aligned}
\]

Kde:
\begin{itemize}
    \item $x$ je vektor rozhodovacích proměnných,
    \item $f(x)$ je účelová funkce,
    \item $g_i(x)$ a $h_j(x)$ jsou omezení,
    \item $X$ určuje typ proměnných, například spojité, celočíselné nebo binární.
\end{itemize}

Typický postup při tvorbě optimalizačního modelu:
\[
\text{reálný problém}
\rightarrow
\text{rozhodovací proměnné}
\rightarrow
\text{účelová funkce}
\rightarrow
\text{omezení}
\rightarrow
\text{výpočet a interpretace řešení}.
\]

Například firma může rozhodovat, kolik vyrobit jednotlivých produktů. Proměnné jsou vyráběná
množství, účelová funkce je zisk a omezení vyjadřují dostupné zdroje, kapacity strojů nebo maximální
poptávku.

\subsection{Lineární programování}

Lineární programování, zkráceně LP, řeší optimalizační úlohy, ve kterých je účelová funkce i všechna
omezení lineární. Standardní maximalizační tvar je
\[
\begin{aligned}
\max_x \quad & c^\top x\\
\text{s.t.}\quad & Ax\le b,\\
& x\ge 0.
\end{aligned}
\]

Kde:
\begin{itemize}
    \item $c$ je vektor koeficientů účelové funkce, například jednotkových zisků,
    \item $A$ je matice technologických nebo strukturálních koeficientů,
    \item $b$ je vektor pravých stran, například dostupných kapacit,
    \item $x$ je vektor rozhodovacích proměnných.
\end{itemize}

Nerovnosti lze převést na rovnosti pomocí doplňkových proměnných:
\[
Ax+s=b,\qquad s\ge 0.
\]
Proměnná $s_i$ říká, kolik zůstalo nevyužité kapacity v $i$-tém omezení. Pokud je $s_i=0$, dané
omezení je aktivní; pokud je $s_i>0$, zdroj není plně využit.

\subsubsection{Geometrická interpretace LP}

Přípustná množina lineárního programu je průnik polorovin, tedy konvexní mnohostěn. Pro LP platí
důležitá vlastnost:

\[
\text{pokud existuje optimální řešení, existuje také v některém vrcholu přípustné množiny.}
\]

Proto není nutné prohledávat všechny body přípustné množiny. Stačí se zaměřit na vrcholy, které
odpovídají tzv. bazickým přípustným řešením.

Možné koncové stavy LP:
\begin{itemize}
    \item jedno optimální řešení,
    \item více optimálních řešení,
    \item neomezené řešení,
    \item žádné přípustné řešení.
\end{itemize}

\subsection{Bazické řešení a bazické přípustné řešení}

Po převodu omezení do tvaru
\[
A'x=b,\qquad x\ge 0
\]
vybíráme určité proměnné jako bazické. Ostatní, nebazické proměnné, nastavíme na nulu.
Pokud výsledné řešení splňuje nezápornost, jde o \textbf{bazické přípustné řešení}.

Je-li v modelu po zavedení slack proměnných $n$ proměnných a $m$ omezení ve tvaru rovností,
pak bázi tvoří $m$ lineárně nezávislých sloupců matice $A'$. Počet možných bází může být velmi
velký:
\[
\binom{n}{m}.
\]
Proto by prosté vyzkoušení všech bazických řešení bylo pro větší úlohy nepraktické.

\subsection{Simplexová metoda}

Simplexová metoda je základní metoda řešení LP. Využívá fakt, že optimum LP leží v některém
vrcholu přípustné množiny. Algoritmus tedy nepohybuje řešením náhodně uvnitř přípustné oblasti,
ale přechází z jednoho bazického přípustného řešení do sousedního, přičemž v každém kroku zlepšuje
hodnotu účelové funkce.

Základní myšlenka:
\[
\text{vrchol}
\rightarrow
\text{sousední lepší vrchol}
\rightarrow
\cdots
\rightarrow
\text{optimální vrchol}.
\]

Pro úlohu
\[
\begin{aligned}
\max_x \quad & c^\top x\\
\text{s.t.}\quad & Ax\le b,\\
& x\ge 0
\end{aligned}
\]
po zavedení slack proměnných vytváříme počáteční simplexovou tabulku:
\[
\left[
\begin{array}{ccc}
A & I & b\\
-c^\top & 0 & 0
\end{array}
\right].
\]

\subsubsection{Postup simplexového algoritmu}

Pro maximalizační úlohu ve standardní tabulkové formě:

\begin{enumerate}
    \item Zkontrolujeme poslední řádek tabulky, tedy řádek účelové funkce.
    \item Pokud v posledním řádku nejsou záporné koeficienty, aktuální řešení je optimální.
    \item Pokud záporné koeficienty existují, vybereme vstupující proměnnou. Obvykle se bere sloupec
    s nejvíce záporným koeficientem.
    \item Tento sloupec se nazývá pivotní sloupec.
    \item Pokud jsou všechny hodnoty v pivotním sloupci nekladné, úloha je neomezená.
    \item Jinak provedeme poměrový test:
    \[
    \frac{b_i}{a_{ik}},\qquad a_{ik}>0.
    \]
    Řádek s nejmenším kladným poměrem určuje vystupující proměnnou.
    \item Prvek v průniku pivotního řádku a pivotního sloupce je pivot.
    \item Pomocí řádkových úprav převedeme pivotní sloupec na jednotkový sloupec.
    \item Dostaneme novou bázi a nové bazické přípustné řešení.
    \item Postup opakujeme.
\end{enumerate}

Intuitivně:
\begin{itemize}
    \item vstupující proměnná je proměnná, kterou se vyplatí začít vyrábět nebo využívat,
    \item vystupující proměnná je proměnná, která se musí z báze odstranit, aby zůstala splněna omezení,
    \item pivotní operace odpovídá přesunu do sousedního vrcholu přípustné množiny.
\end{itemize}

\subsubsection{Jednoduchý příklad LP}

Firma vyrábí produkty $A$ a $B$:
\[
\begin{aligned}
\max \quad & 2x_1+x_2\\
\text{s.t.}\quad & x_1+2x_2\le 8,\\
& x_1\le 6,\\
& x_1,x_2\ge 0.
\end{aligned}
\]

Po zavedení slack proměnných:
\[
\begin{aligned}
x_1+2x_2+x_3&=8,\\
x_1+x_4&=6.
\end{aligned}
\]

Počáteční řešení je
\[
x=(0,0,8,6),
\]
tedy firma zatím nevyrábí nic, ale má nevyužité kapacity. Simplexová metoda postupně zjišťuje,
který produkt se vyplatí zařadit do báze, a přepočítává řešení, dokud není dosaženo optima.

\subsubsection{Poznámky ke simplexové metodě}

\begin{itemize}
    \item V nejhorším případě má simplexová metoda exponenciální složitost.
    \item V praxi je však často velmi rychlá a efektivní.
    \item Pokud není zřejmé počáteční bazické přípustné řešení, používá se dvoufázová simplexová metoda.
    \item Revidovaná simplexová metoda neukládá celou tabulku, ale pracuje efektivněji s maticovou reprezentací báze.
    \item Pro velké úlohy se používají i metody vnitřních bodů, které se nepohybují po hranici přípustné množiny,
    ale uvnitř ní.
\end{itemize}

\subsection{Dualita a ekonomická interpretace}

K primalní úloze
\[
\begin{aligned}
(P)\quad \max_x \quad & c^\top x\\
\text{s.t.}\quad & Ax\le b,\\
& x\ge 0
\end{aligned}
\]
existuje duální úloha
\[
\begin{aligned}
(D)\quad \min_y \quad & b^\top y\\
\text{s.t.}\quad & A^\top y\ge c,\\
& y\ge 0.
\end{aligned}
\]

Duální proměnné $y_i$ mají ekonomickou interpretaci jako \textbf{stínové ceny}. Stínová cena ukazuje,
o kolik by se změnila optimální hodnota účelové funkce, pokud bychom zvýšili pravou stranu daného
omezení o jednu jednotku:
\[
y_i^* \approx \frac{\partial z^*}{\partial b_i}.
\]

Například pokud je omezením počet hodin práce a stínová cena je $5$, znamená to, že jedna dodatečná
hodina práce zvýší optimální zisk přibližně o $5$ jednotek. Tato interpretace platí jen v určitém
intervalu změn, protože při větší změně se může změnit optimální báze.

\subsubsection{Slabá a silná dualita}

Slabá dualita říká, že pro každé přípustné řešení primalu a duálu platí
\[
c^\top x \le b^\top y.
\]
Duální úloha tedy poskytuje horní mez pro primalní maximalizační úlohu.

Silná dualita říká, že pokud existuje optimální řešení, pak optimální hodnoty primalu a duálu jsou stejné:
\[
c^\top x^* = b^\top y^*.
\]

\subsection{Typické úlohy lineárního programování}

\subsubsection{Výrobní plánování}

Firma vyrábí produkty $j=1,\dots,n$ a využívá zdroje $i=1,\dots,m$:
\[
\begin{aligned}
\max_x \quad & \sum_{j=1}^n p_jx_j\\
\text{s.t.}\quad & \sum_{j=1}^n a_{ij}x_j\le b_i,\quad i=1,\dots,m,\\
& x_j\ge 0.
\end{aligned}
\]

Kde:
\begin{itemize}
    \item $x_j$ je množství produktu $j$,
    \item $p_j$ je jednotkový zisk produktu $j$,
    \item $a_{ij}$ je spotřeba zdroje $i$ na jednotku produktu $j$,
    \item $b_i$ je dostupná kapacita zdroje $i$.
\end{itemize}

\subsubsection{Dopravní problém}

Cílem je minimalizovat náklady přepravy ze skladů k odběratelům:
\[
\begin{aligned}
\min_x \quad & \sum_i\sum_j c_{ij}x_{ij}\\
\text{s.t.}\quad & \sum_j x_{ij}\le s_i,\quad \forall i,\\
& \sum_i x_{ij}\ge d_j,\quad \forall j,\\
& x_{ij}\ge 0.
\end{aligned}
\]

\subsubsection{Směšovací a dietní úlohy}

Cílem je najít nejlevnější kombinaci vstupů při splnění požadavků na složení:
\[
\begin{aligned}
\min_x \quad & c^\top x\\
\text{s.t.}\quad & Ax\ge q,\\
& x\ge 0.
\end{aligned}
\]

Použití:
\begin{itemize}
    \item krmné směsi,
    \item energetický mix,
    \item investiční portfolio,
    \item plánování nákupu surovin.
\end{itemize}

\subsection{Celočíselné a smíšené celočíselné programování}

Celočíselné lineární programování, zkráceně ILP, má lineární účelovou funkci a lineární omezení,
ale některé proměnné musí být celočíselné:
\[
\begin{aligned}
\max_x \quad & c^\top x\\
\text{s.t.}\quad & Ax\le b,\\
& x\in \mathbb{Z}^n.
\end{aligned}
\]

Pokud jsou celočíselné pouze některé proměnné a ostatní jsou spojité, jde o smíšené celočíselné
lineární programování, tedy MILP:
\[
\begin{aligned}
\max_{x,y}\quad & c^\top x+d^\top y\\
\text{s.t.}\quad & Ax+By\le b,\\
& x\ge 0,\\
& y\in \mathbb{Z}^p.
\end{aligned}
\]

Zvlášť důležité jsou binární proměnné:
\[
y_j\in\{0,1\}.
\]
Ty reprezentují rozhodnutí typu ano/ne:
\[
y_j=
\begin{cases}
1, & \text{varianta se zvolí,}\\
0, & \text{varianta se nezvolí.}
\end{cases}
\]

\subsubsection{Proč nestačí zaokrouhlit LP řešení}

U ILP nelze obecně vyřešit LP relaxaci a výsledek pouze zaokrouhlit. Důvody:
\begin{itemize}
    \item zaokrouhlené řešení může porušit omezení,
    \item zaokrouhlené řešení může být přípustné, ale daleko od optima,
    \item optimum LP relaxace často leží ve vrcholu, který není celočíselný.
\end{itemize}

Příklad: pokud omezení říká
\[
x_1+x_2\le 1
\]
a LP řešení je
\[
x_1=x_2=0.6,
\]
pak zaokrouhlení na
\[
x_1=x_2=1
\]
vede k porušení omezení.

\subsection{LP relaxace a integrality gap}

LP relaxace vznikne odstraněním celočíselných podmínek. Například:
\[
x_i\in\{0,1\}
\quad\Rightarrow\quad
0\le x_i\le 1.
\]

Pro maximalizační úlohu platí:
\[
z^*_{LP}\ge z^*_{ILP},
\]
protože LP relaxace má větší přípustnou množinu.

Pro minimalizační úlohu platí:
\[
z^*_{LP}\le z^*_{ILP}.
\]

Při řešení ILP se pracuje s mezemi:
\begin{itemize}
    \item u minimalizace poskytuje LP relaxace dolní mez,
    \item každé přípustné celočíselné řešení poskytuje horní mez,
    \item u maximalizace je interpretace opačná.
\end{itemize}

Relativní mezera se často zapisuje jako
\[
gap=
\frac{z_{UB}-z_{LB}}{|z_{UB}|},
\]
kde $z_{UB}$ je horní mez a $z_{LB}$ dolní mez. Pokud je mezera malá, lze řešení prakticky považovat
za dostatečně dobré.

\subsection{Branch-and-bound}

Branch-and-bound je základní přesná metoda pro řešení ILP/MILP. Kombinuje řešení LP relaxací
s postupným dělením problému na menší podproblémy.

Postup pro maximalizační úlohu:
\begin{enumerate}
    \item Vyřešíme LP relaxaci původní úlohy.
    \item Pokud je řešení celočíselné, je kandidátem na optimum.
    \item Pokud řešení není celočíselné, vybereme neceločíselnou proměnnou $x_i^*$.
    \item Provedeme větvení:
    \[
    x_i\le \lfloor x_i^*\rfloor,
    \qquad
    x_i\ge \lceil x_i^*\rceil.
    \]
    \item Vzniknou dva nové podproblémy.
    \item Každý podproblém řešíme jako LP relaxaci.
    \item Větev zahodíme, pokud:
    \begin{itemize}
        \item je nepřípustná,
        \item její LP relaxace nemůže překonat nejlepší dosud nalezené celočíselné řešení,
        \item její LP řešení je celočíselné a bylo vyhodnoceno.
    \end{itemize}
\end{enumerate}

Důležitý pojem je \textbf{incumbent}, tedy nejlepší dosud nalezené celočíselné řešení. Branch-and-bound
systematicky zmenšuje prostor řešení, dokud neprokáže optimálnost incumbentu nebo nedosáhne
požadované mezery.

\subsection{Cutting planes}

Metoda cutting planes řeší ILP pomocí dodatečných lineárních omezení, tzv. řezů. Řez je omezení,
které odstraní část přípustné množiny LP relaxace, ale neodstraní žádné přípustné celočíselné řešení.

Formálně: nechť
\[
P_R=\{x\in\mathbb{R}^n:Ax\le b\}
\]
je relaxovaná přípustná množina a
\[
P_I=\{x\in\mathbb{Z}^n:Ax\le b\}
\]
je množina celočíselných řešení. Řez má tvar
\[
\alpha^\top x\le \beta
\]
a musí splňovat:
\[
\alpha^\top x\le \beta \quad \text{pro všechna } x\in P_I,
\]
ale současně má být porušen aktuálním neceločíselným LP řešením $x^R$:
\[
\alpha^\top x^R> \beta.
\]

Tím se aktuální neceločíselné řešení odřízne, ale všechna celočíselná řešení zůstanou zachována.

\subsubsection{Algoritmus cutting plane}

\begin{enumerate}
    \item Vyřešíme LP relaxaci.
    \item Pokud je optimální řešení celočíselné, končíme.
    \item Pokud je řešení neceločíselné, najdeme platný řez, který jej odřízne.
    \item Řez přidáme do modelu jako nové omezení.
    \item Znovu řešíme LP relaxaci.
\end{enumerate}

Intuice:
\[
\text{velká LP oblast}
\rightarrow
\text{přidání řezů}
\rightarrow
\text{menší LP oblast}
\rightarrow
\text{vrcholy blíže celočíselným řešením}.
\]

Po konečném počtu vhodných řezů se může stát, že všechny vrcholy relaxované množiny jsou
celočíselné. Potom řešení LP poskytne i řešení ILP.

\subsubsection{Gomoryho řezy}

Gomoryho řezy jsou jedním ze známých způsobů generování řezů. Vycházejí ze simplexové tabulky.
Pokud má některá bazická proměnná neceločíselnou hodnotu, lze z jejího řádku vytvořit omezení,
které aktuální neceločíselné řešení odřízne.

Uvažujme řádek simplexové tabulky:
\[
x_{B_i}=b_i-\sum_{j\in N}a_{ij}x_j,
\]
kde $b_i$ není celé číslo. Označme zlomkovou část čísla jako
\[
\{r\}=r-\lfloor r\rfloor.
\]
Potom lze zapsat Gomoryho řez například ve tvaru
\[
\sum_{j\in N}\{a_{ij}\}x_j \ge \{b_i\}.
\]

Tento řez:
\begin{itemize}
    \item je platný pro všechna celočíselná řešení,
    \item odřízne aktuální zlomkové řešení LP relaxace,
    \item lze relativně efektivně generovat ze simplexové tabulky.
\end{itemize}

V praxi se samotné cutting planes často kombinují s branch-and-bound. Tato kombinace se nazývá
\textbf{branch-and-cut}. Moderní MILP řešiče používají právě tento princip: v uzlech větvicího stromu
řeší LP relaxace, přidávají řezy a teprve potom větví dál.

\subsubsection{Řezy v TSP}

U problému obchodního cestujícího se řezy používají například k eliminaci podcyklů. Pokud $S$ je
podmnožina měst, pak nesmí vzniknout samostatný cyklus pouze uvnitř $S$:
\[
\sum_{\substack{i,j\in S\\ i<j}}x_{ij}\le |S|-1.
\]

Tato omezení se nazývají subtour elimination constraints. Protože jich je exponenciálně mnoho,
nepřidávají se všechna najednou. Řešič nejprve vyřeší relaxovanou úlohu a potom přidává pouze ta
omezení, která jsou aktuálním řešením porušena.

\subsection{Typické úlohy MILP}

\subsubsection{Batohový problém}

Máme položky s hodnotou $p_i$ a váhou $w_i$. Chceme vybrat nejhodnotnější kombinaci, která se
vejde do kapacity $W$:
\[
\begin{aligned}
\max_x \quad & \sum_{i=1}^n p_ix_i\\
\text{s.t.}\quad & \sum_{i=1}^n w_ix_i\le W,\\
& x_i\in\{0,1\}.
\end{aligned}
\]

Proměnná $x_i=1$ znamená, že položku vybereme. Ekonomická interpretace může být výběr projektů
při omezeném rozpočtu.

\subsubsection{Lokalizační úloha se fixními náklady}

Rozhodujeme, která zařízení otevřít a jak k nim přiřadit zákazníky:
\[
\begin{aligned}
\min_{x,y}\quad
& \sum_j f_jy_j+\sum_i\sum_j c_{ij}x_{ij}\\
\text{s.t.}\quad
& \sum_j x_{ij}=1,\quad \forall i,\\
& x_{ij}\le y_j,\quad \forall i,j,\\
& y_j\in\{0,1\},\\
& x_{ij}\ge 0.
\end{aligned}
\]

Kde:
\begin{itemize}
    \item $y_j=1$ znamená, že zařízení $j$ otevřeme,
    \item $x_{ij}$ vyjadřuje přiřazení zákazníka $i$ k zařízení $j$,
    \item $f_j$ jsou fixní náklady otevření zařízení,
    \item $c_{ij}$ jsou náklady obsluhy zákazníka.
\end{itemize}

Omezení
\[
x_{ij}\le y_j
\]
zajišťuje, že zákazníka nelze přiřadit k neotevřenému zařízení.

\subsubsection{Modelování logických podmínek pomocí Big-M}

Binární proměnné umožňují modelovat logická rozhodnutí. Například pokud výroba $x$ může probíhat
pouze tehdy, pokud je aktivována varianta $y$, lze použít:
\[
x\le My,\qquad y\in\{0,1\}.
\]

Pokud $y=0$, potom $x\le 0$ a výroba neprobíhá. Pokud $y=1$, potom $x\le M$ a výroba je povolena.
Konstanta $M$ má být dostatečně velká, ale ne zbytečně obrovská, protože příliš velké $M$ zhoršuje
numerickou stabilitu a kvalitu LP relaxace.

\subsection{Kombinatorická optimalizace}

Kombinatorická optimalizace hledá nejlepší objekt z konečné množiny objektů. Typicky jde o cesty,
cykly, množiny hran, rozvrhy nebo přiřazení. Často se používá v grafových úlohách.

\subsubsection{Nejkratší cesta}

Máme graf $G=(V,A)$ a délky hran $c_{ij}$. Chceme najít nejkratší cestu ze zdroje $s$ do cíle $t$:
\[
\min \sum_{(i,j)\in A} c_{ij}x_{ij}.
\]

Toková formulace:
\[
\sum_{j:(i,j)\in A}x_{ij}
-
\sum_{j:(j,i)\in A}x_{ji}
=
\begin{cases}
1, & i=s,\\
-1, & i=t,\\
0, & \text{jinak.}
\end{cases}
\]

Použití:
\begin{itemize}
    \item logistika,
    \item navigace,
    \item routing,
    \item plánování toku materiálu.
\end{itemize}

\subsubsection{Minimální kostra grafu}

Máme neorientovaný souvislý graf $G=(V,E)$ a chceme propojit všechny uzly s minimální celkovou
vahou bez vytvoření cyklu:
\[
\min \sum_{e\in E} w_ex_e.
\]

Základní podmínky:
\[
\sum_{e\in E}x_e=|V|-1,
\]
\[
\sum_{e\in E(S)}x_e\le |S|-1,\qquad \forall S\subseteq V.
\]

Interpretace:
\begin{itemize}
    \item návrh nejlevnější sítě,
    \item propojení poboček,
    \item návrh komunikační nebo distribuční infrastruktury.
\end{itemize}

\subsubsection{Problém obchodního cestujícího}

Problém obchodního cestujícího, TSP, hledá nejkratší cyklus, který navštíví každé město právě jednou
a vrátí se do výchozího města.

Pro symetrický TSP:
\[
\begin{aligned}
\min_x \quad & \sum_{i<j}d_{ij}x_{ij}\\
\text{s.t.}\quad
& \sum_{j<i}x_{ji}+\sum_{j>i}x_{ij}=2,\quad \forall i\in V,\\
& \sum_{\substack{i,j\in S\\ i<j}}x_{ij}\le |S|-1,\quad \forall S\subset V,\; 3\le |S|\le |V|-1,\\
& x_{ij}\in\{0,1\}.
\end{aligned}
\]

První skupina omezení říká, že do každého města vstupují právě dvě hrany. Druhá skupina omezení
eliminuje podcykly.

TSP je důležitý, protože je jednoduchý na formulaci, ale výpočetně náročný. Slouží jako typický příklad
NP-těžké úlohy a testovací problém pro obecné optimalizační metody.

\subsection{Heuristiky a metaheuristiky}

U velkých ILP/MILP a kombinatorických úloh může být přesné nalezení optima příliš časově náročné.
Proto se používají heuristiky a metaheuristiky.

\subsubsection{Heuristika}

Heuristika je postup, který hledá dobré přípustné řešení, ale obvykle negarantuje nalezení globálního
optima. Může být:
\begin{itemize}
    \item konstruktivní -- staví řešení od začátku,
    \item zlepšovací -- začíná s existujícím řešením a snaží se ho vylepšit.
\end{itemize}

U minimalizační úlohy poskytuje každé přípustné heuristické řešení horní mez:
\[
z_{UB}=z(x^{heur}).
\]

\subsubsection{Metaheuristika}

Metaheuristika je obecnější rámec, který není vázán jen na jeden konkrétní typ problému.
Příklady:
\begin{itemize}
    \item lokální hledání,
    \item simulované žíhání,
    \item genetické algoritmy,
    \item ant colony optimization.
\end{itemize}

Výhoda metaheuristik:
\begin{itemize}
    \item použitelnost pro mnoho typů úloh,
    \item schopnost najít dobré řešení u velkých problémů,
    \item možnost pracovat i tam, kde přesné metody selhávají časově.
\end{itemize}

Nevýhoda:
\begin{itemize}
    \item bez záruky globálního optima,
    \item výsledek závisí na parametrech a počátečním řešení,
    \item často je vhodné provádět více běhů algoritmu.
\end{itemize}

\subsection{Lokální hledání}

Lokální hledání vychází z aktuálního řešení a zkoumá jeho okolí:
\[
x^{(t+1)}\in N(x^{(t)}).
\]

Okolí $N(x)$ je množina řešení, která se od aktuálního řešení liší jen malou změnou.

Příklady okolí:
\begin{itemize}
    \item u TSP prohození dvou hran,
    \item u binárního programu změna jedné proměnné z $0$ na $1$ nebo opačně,
    \item u rozvrhování prohození dvou úloh.
\end{itemize}

Postup:
\begin{enumerate}
    \item začneme z nějakého přípustného řešení,
    \item najdeme lepší řešení v jeho okolí,
    \item přesuneme se do něj,
    \item opakujeme, dokud nelze najít zlepšení.
\end{enumerate}

Problémem je lokální optimum. Lokální optimum je řešení, které je nejlepší ve svém okolí, ale nemusí být
globálně nejlepší.

\subsection{Simulované žíhání}

Simulované žíhání rozšiřuje lokální hledání tím, že občas přijme i horší řešení. To umožňuje uniknout
z lokálního optima.

Pro minimalizační úlohu označme:
\[
\Delta=f(x')-f(x).
\]

Nové řešení $x'$ přijmeme, pokud:
\[
\Delta<0,
\]
nebo s pravděpodobností
\[
P=\exp\left(-\frac{\Delta}{T}\right),
\]
kde $T$ je teplota.

Interpretace:
\begin{itemize}
    \item při vysoké teplotě algoritmus více experimentuje,
    \item při nízké teplotě se chová podobně jako lokální hledání,
    \item postupné snižování teploty se nazývá cooling schedule.
\end{itemize}

\subsection{Genetický algoritmus}

Genetický algoritmus je inspirován biologickou evolucí. Nepracuje s jedním řešením, ale s celou
populací řešení.

Pojmy:
\begin{itemize}
    \item jedinec = jedno kandidátní řešení,
    \item populace = množina kandidátních řešení,
    \item fitness = kvalita řešení,
    \item generace = jedna iterace vývoje populace.
\end{itemize}

Základní kroky:
\begin{enumerate}
    \item inicializace počáteční populace,
    \item výpočet fitness,
    \item výběr rodičů,
    \item křížení,
    \item mutace,
    \item vytvoření nové generace,
    \item elitismus -- zachování nejlepších řešení.
\end{enumerate}

Mutace udržuje diverzitu populace. Křížení kombinuje vlastnosti dvou dobrých řešení.
Elitismus zajišťuje, že se nejlepší dosud nalezené řešení neztratí.

\subsection{Ant Colony Optimization}

Ant colony optimization je metaheuristika inspirovaná chováním mravenců. Mravenci při hledání potravy
zanechávají feromonové stopy. Cesty, které vedou k dobrým řešením, jsou postupně posilovány.

U TSP si lze představit, že každý umělý mravenec postupně konstruuje trasu. Pravděpodobnost přechodu
z města $i$ do města $j$ může být úměrná:
\[
p_{ij}\propto \tau_{ij}^{\alpha}\eta_{ij}^{\beta},
\]
kde:
\begin{itemize}
    \item $\tau_{ij}$ je množství feromonu na hraně,
    \item $\eta_{ij}$ je heuristická informace, například $\eta_{ij}=1/d_{ij}$,
    \item $\alpha$ a $\beta$ určují váhu feromonu a heuristiky.
\end{itemize}

Aktualizace feromonu:
\[
\tau_{ij}\leftarrow (1-\rho)\tau_{ij}+\Delta\tau_{ij},
\]
kde $\rho$ je míra vypařování feromonu. Vypařování brání tomu, aby algoritmus příliš brzy zafixoval
špatnou cestu.

\subsection{Penalizace omezení}

Pokud je obtížné pracovat přímo s omezeními, lze je zahrnout do účelové funkce pomocí penalizace.
Pro problém
\[
\min f(x)
\quad \text{s.t.}\quad
g_i(x)\le 0
\]
lze vytvořit penalizovanou funkci:
\[
F(x)=f(x)+c\sum_{i=1}^m \max\{0,g_i(x)\}^2.
\]

Pokud je omezení splněno, člen $\max\{0,g_i(x)\}$ je nulový. Pokud je omezení porušeno, účelová
funkce se zhorší. Parametr $c$ určuje sílu penalizace.

\subsection{Gradientní metody}

U spojitých nelineárních úloh lze použít numerické metody. Gradientní sestup aktualizuje řešení podle
vzorce:
\[
x^{new}=x^{old}-\eta\nabla f(x^{old}),
\]
kde:
\begin{itemize}
    \item $\nabla f(x)$ je gradient,
    \item $\eta$ je velikost kroku.
\end{itemize}

Gradient ukazuje směr největšího růstu funkce, proto při minimalizaci postupujeme opačným směrem.
Příliš malý krok vede k pomalé konvergenci, příliš velký krok může způsobit oscilace nebo divergenci.

Newtonova metoda navíc využívá Hessovu matici:
\[
H(f)=
\left(
\frac{\partial^2 f}{\partial x_i\partial x_j}
\right).
\]
Díky informaci o zakřivení může konvergovat rychleji, ale je výpočetně náročnější.

\subsection{Shrnutí ke zkoušce}

\begin{itemize}
    \item Optimalizační model má rozhodovací proměnné, účelovou funkci a omezení.
    \item LP má lineární účelovou funkci a lineární omezení.
    \item Přípustná množina LP je konvexní a optimum, pokud existuje, leží ve vrcholu.
    \item Simplexová metoda postupuje mezi bazickými přípustnými řešeními a zlepšuje účelovou funkci.
    \item Duální proměnné LP lze interpretovat jako stínové ceny zdrojů.
    \item ILP/MILP přidává celočíselné a binární proměnné, čímž umožňuje modelovat výběr, fixní náklady,
    logické podmínky a kombinatorické struktury.
    \item LP relaxace odstraňuje celočíselnost a poskytuje mez pro původní celočíselnou úlohu.
    \item Branch-and-bound řeší ILP pomocí LP relaxací, větvení a ořezávání neperspektivních větví.
    \item Cutting planes přidávají platná omezení, která odříznou neceločíselné LP řešení, ale zachovají všechna
    celočíselná řešení.
    \item Moderní řešiče často kombinují branch-and-bound a cutting planes do metody branch-and-cut.
    \item Typické úlohy: výrobní plánování, dopravní problém, batoh, lokalizace zařízení, nejkratší cesta,
    minimální kostra a TSP.
    \item Heuristiky a metaheuristiky hledají dobrá řešení bez záruky optima, ale jsou prakticky důležité u velkých
    a výpočetně těžkých úloh.
\end{itemize}
